% !TEX program = lualatex
\documentclass[12pt, a4paper]{article}
\usepackage{master}

\DeclareWorkGR{Cat}{범주론}{Κατηγορίαι}{Categoriae}
\DeclareWorkGR{DI}{명제론}{Περὶ ἑρμηνείας}{De interpretatione}

\fancyhead[L]{\footnotesize\color{midgray}오르가논 강독}
\fancyhead[R]{\footnotesize\color{midgray}특별 세션 β}

\begin{document}
\thispagestyle{firstpage}

\weeksection{특별 세션 β}{범주론·이사고게 종합 II}{수용사 --- 세 문명권의 천 년}

\subsection{이슬람 세계의 수용: 범주론은 논리학에 속하는가}

\subsubsection{알렉산드리아에서 바그다드까지 --- 시리아어 다리}

이사고게-범주론 체계가 이슬람 세계로 전해진 경로는 시리아 기독교인들을 통해서였습니다. 5--6세기 알렉산드리아 학파는 오르가논의 축약된 교육 과정을 사용했는데, 이사고게 + 범주론 + 명제론 + 분석론 전서 I.1--7까지만 가르치고, 양상 삼단논법(I.8 이후)은 `읽지 않는 부분'으로 남겨두었습니다.\footnote{SEP, ``al-Farabi's Philosophy of Logic and Language'' (Hodges, 2023).} 6세기 레사이나의 세르기우스(Sergius of Reshaina, d.\ 536)가 암모니오스의 제자로서 범주론을 시리아어로 번역한 것이 전승의 출발점이며, 9세기 후나인 이븐 이스하크(\d{H}unayn ibn Is\d{h}\=aq)의 번역으로 완전한 아랍어 오르가논이 완성되었습니다.

\subsubsection{알킨디와 알파라비 --- 1차 개념과 2차 개념}

알킨디(al-Kind\=\i, c.\ 800--870)는 양과 질에 특별한 지위를 부여하여 나머지 범주들을 이 둘과 실체의 조합으로 설명하는 환원론적 독해를 시도했습니다.\footnote{Ahmad Ighbariah, ``Al-Kind\=\i{} and the Reception of Aristotle's Categories,'' \gri{Arabic Sciences and Philosophy} 22, no.\ 1 (2012).} 4주차에서 다룬 질 범주의 통일성 문제를 떠올려 봅시다 --- 알킨디는 퍼스의 `기괴한 오합지졸'과 정반대로, 질을 범주 체계의 기둥 가운데 하나로 격상시켰습니다.

알파라비(al-F\=ar\=ab\=\i, c.\ 870--950), `제2의 스승'은 결정적인 구별을 도입했습니다 --- 1차 개념(ma\textquotesingle q\=ul\=at \=ul\=a)과 2차 개념(ma\textquotesingle q\=ul\=at th\=aniya).\footnote{SEP, ``al-Farabi's Philosophy of Logic and Language'' (Hodges, 2023).} 1차 개념은 감각 경험에서 직접 끌어온 것(`사람', `흰')이고, 2차 개념은 개념에 대한 개념(`류', `종', `더 포괄적인 것')입니다. 이사고게의 다섯 술어가능자는 2차 개념이고, 범주론의 10범주는 1차 개념입니다. 논리학이 2차 개념을 연구하는 학문이라면, 1차 개념을 분류하는 범주론은 엄밀한 의미에서 논리학에 속하지 않게 됩니다.

\subsubsection{아비첸나의 혁명 --- 범주론을 논리학에서 추방하다}

아비첸나(Ibn S\=\i n\=a, 980--1037)는 알파라비의 구별에서 급진적 결론을 이끌어냈습니다.\footnote{SEP, ``Ibn Sina's Logic'' (Strobino).} 논리학의 주제는 2차 지향 --- `주어', `술어', `전제', `삼단논법' 같은 개념에 관한 개념들 --- 입니다. 범주론은 1차 지향을 분류하므로 형이상학에 속합니다. 아비첸나는 범주론을 `페리파토스 학파의 관습에 대한 경의'로만 유지했습니다.

아비첸나의 종차 비판도 우리의 강독과 직접 연결됩니다. 7주차에서 포르피리오스가 종차를 `종에 있어서 서로 다른 여럿에 대해 술어되는 것'으로 정의한 것을 보았습니다. 아비첸나는 이를 거부하고, 종차를 `류 아래에서 종에 대해 그것이 무엇인지를 본질적으로 대답하는 단순한 보편자'로 재정의했습니다.\footnote{Cristina Cerami, ``Avicenna against Porphyry's Definition of Differentia Specifica,'' \gri{Documenti e studi sulla tradizione filosofica medievale} 26 (2015): 129--183.}

아비첸나의 보편자론도 독자적입니다. `말성(馬性)은 오직 말성일 뿐이다' --- 본질 자체는 보편적이지도 개별적이지도 않습니다. 이슬람 논리학자들은 이후 보편자를 자연적 보편자(본질 자체), 논리적 보편자(술어 가능성), 지성적 보편자(추상된 개념)로 삼분했습니다.

\subsubsection{아베로에스 --- 정통으로의 복귀}

아베로에스(Ibn Rushd, 1126--1198)는 아비첸나에 맞서 아리스토텔레스 정통으로의 복귀를 시도하면서, 흥미롭게도 이사고게를 `논리학의 입문도 아니고 일부도 아니다'라는 이유로 오르가논에서 분리했습니다.

{\footnotesize
\begin{concepttable}{|l|c|c|c|}
\hline
\rowcolor{tableheadblue}
{\color{h1navy}\bfseries 학자} &
{\color{h1navy}\bfseries 범주론의 주제} &
{\color{h1navy}\bfseries 논리학?} &
{\color{h1navy}\bfseries 이사고게 지위} \\ \hline
알렉산드로스 (c.\ 200) & 사물 & 형이상학 & (이전) \\ \hline
포르피리오스 (c.\ 280) & 의미하는 표현 & 논리학 & 입문 \\ \hline
알파라비 (c.\ 940) & 1차 개념 & 논리학(교육적) & 입문 \\ \hline
아비첸나 (c.\ 1020) & 1차 지향 & 형이상학 & 입문 유지 \\ \hline
아베로에스 (c.\ 1180) & 사물 & 논리학(복귀) & 분리 \\ \hline
\end{concepttable}
}

\subsection{비잔틴 세계: 고대 필사본의 보존자들}

\subsubsection{비잔틴 철학 교육에서의 오르가논}

비잔틴 세계(330--1453)는 그리스어 원문을 직접 읽을 수 있었던 유일한 중세 문명입니다. 오르가논의 비잔틴 필사본은 100부 이상 현존하며, 성경 다음으로 많은 숫자입니다.\footnote{Katerina Ierodiakonou, ``Logic in Byzantium,'' in \gri{The Cambridge Intellectual History of Byzantium}, ed.\ A.\ Kaldellis and N.\ Siniossoglou (Cambridge: Cambridge University Press, 2017).}

\subsubsection{프셀로스, 이탈로스, 프로드로모스}

미하엘 프셀로스(Michael Psellos, 1018--1078)는 류와 종을 이해하는 세 가지 방식(ante rem, in re, post rem)을 암시하여 `온건한 실재론'으로 분류됩니다.\footnote{SEP, ``Byzantine Philosophy'' (Ierodiakonou); K.\ Benakis, ``Michael Psellos,'' in \gri{Philosophie: Geschichtliche Einführungen}, vol.\ 6 (Basel, 1982).} 그의 제자 요안네스 이탈로스(John Italos, c.\ 1025--1082 이후)는 안나 콤네네에 따르면 `아리스토텔레스를 당대 누구보다 잘 해석했다'고 하나, 1082년 교회에 의해 정죄되었습니다 --- 철학적 논리를 신학에 지나치게 적용했다는 이유로. 이것은 1277년 파리 금지령보다 200년 앞선 범주론-신학 충돌 사례입니다.

테오도로스 프로드로모스(c.\ 1100--1170)는 크세네데모스(\gr{Ξενέδημος})라는 대화편에서 포르피리오스의 다섯 술어가능자를 논의했고, `크다'와 `작다'에 관하여에서 아리스토텔레스가 이것들을 양이 아니라 관계로 분류한 것을 비판했습니다 --- 4주차의 범주 간 교차 문제의 또 다른 사례입니다.

\subsubsection{필사본의 여행 --- 비잔틴에서 라틴 서방으로}

비잔틴의 가장 중요한 공헌은 그리스어 필사본의 보존과 전달입니다. 12세기에 베네치아의 야코부스가 콘스탄티노폴리스의 필사본에서 분석론 후서를 번역하면서 신논리학(logica nova)이 탄생합니다. 15세기 스콜라리오스(Scholarios)는 아퀴나스의 영향을 받은 이사고게-범주론 주석서를 남겨, 역방향(라틴 → 그리스) 전승의 사례를 보여줍니다.

\subsection{라틴 서방: 보에티우스의 다리에서 보편자 전쟁까지}

\subsubsection{보에티우스 --- 6세기 동안 유일한 다리}

보에티우스(c.\ 480--524)는 이사고게에 대한 제2주석서에서, 원문의 불과 한 쪽 남짓(포르피리오스의 세 물음)에 대해 제1권 전체를 할애합니다.\footnote{Paul Vincent Spade, \gri{Five Texts on the Mediaeval Problem of Universals} (Indianapolis: Hackett, 1994), 20--25.} 그의 해법은 추상론이었지만, 플라톤과 아리스토텔레스 사이에서 결정을 유보했습니다. 이 이중적 유보가 중세 보편자 논쟁의 구조적 틀을 만들었습니다.

\subsubsection{구논리학에서 신논리학으로}

6세기부터 12세기 중반까지 접근 가능했던 구논리학(logica vetus)은 이사고게, 범주론, 명제론, 보에티우스 논고뿐이었습니다.\footnote{L.\ M.\ De Rijk, \gri{Logica Modernorum} (Assen: Van Gorcum, 1962).} 나머지 오르가논은 12세기 후반에 신논리학(logica nova)으로 도착합니다.

{\footnotesize
\begin{concepttable}{|l|l|l|}
\hline
\rowcolor{tableheadblue}
{\color{h1navy}\bfseries 시기} &
{\color{h1navy}\bfseries 접근 가능한 텍스트} &
{\color{h1navy}\bfseries 명칭} \\ \hline
6--12세기 전반 & 이사고게 + 범주론 + 명제론 + 보에티우스 & 구논리학 \\ \hline
12세기 후반-- & + 분석론 + 토피카 + 소피스트적 논박 & 신논리학 \\ \hline
13세기-- & + 대치·귀결·의무·불용해 등 & 근대 논리학 \\ \hline
\end{concepttable}
}

9세기 서방의 학생이 논리학을 배운다고 상상해 봅시다. 접할 수 있는 것은 이사고게, 범주론, 명제론뿐입니다. 삼단논법은 보에티우스의 축약본으로만 접합니다. 정의와 논증의 관계(분석론 후서), 변증적 논쟁(토피카)은 아예 알 수 없습니다. 12세기에 나머지 오르가논이 도착했을 때의 지적 충격은 --- 비유하자면 --- 갈릴레오 이전의 천문학자에게 갑자기 망원경이 주어진 것과 비교할 수 있을 것입니다.

\subsubsection{보편자 논쟁의 주요 입장들}

포르피리오스의 세 물음이 촉발한 논쟁의 주요 입장들을 정리합시다.\footnote{SEP, ``The Medieval Problem of Universals'' (Klima, rev.\ 2022).}

로스켈리누스(c.\ 1050--1125)는 보편자가 `목소리의 내뿜음'(flatus vocis)에 불과하다고 주장했으며, 소아송 공의회(1092)에서 정죄되었습니다. 아벨라르(1079--1142)는 양쪽 극단을 비판하고, 보편자를 의미 있는 단어(sermo)로 보면서 `상태'(status) 개념을 도입했습니다. 아퀴나스(1225--1274)는 ante rem / in re / post rem 삼중 구조를 구축했고, 스코투스(c.\ 1266--1308)는 `공통 본성' + `이것임'(haecceitas)으로, 오캄(c.\ 1287--1347)은 마음의 개념(intentio animae)으로 각각 해법을 제시했습니다.

여기서 사고 연습을 해봅시다. 아비첸나의 삼분법(자연적/논리적/지성적 보편자)과 라틴 삼분법(ante rem/in re/post rem)은 정확히 대응하지 않습니다. 전자는 보편자의 `무엇임'을, 후자는 `어디에 있음'을 묻습니다.

\subsubsection{신학 속의 범주론 --- 삼위일체와 화체}

범주론의 실체/부수성 구별은 기독교 신학에서 두 가지 핵심적 역할을 했습니다.

첫째, 삼위일체론. 카파도키아 교부들이 실체(\gr{οὐσία})와 위격(\gr{ὑπόστασις})을 구별하여 `하나의 실체, 세 위격' 공식을 확립했습니다.\footnote{Basil, Epistle 214.4.} 보에티우스는 삼위일체론(De Trinitate)에서 범주론의 10범주를 신에게 적용하여 관계 범주가 핵심이라고 주장했습니다: `실체의 범주는 단일성을 보존하고, 관계의 범주는 삼위를 산출한다'(substantia continet unitatem, relatio multiplicat trinitatem).\footnote{Boethius, \gri{De Trinitate} VI.}

둘째, 화체론. 제4차 라테란 공의회(1215)는 빵의 실체가 변환되되 부수성은 남는다고 선언했습니다. 아퀴나스(신학대전 III, q.\ 77)는 신적 전능으로 실체 없는 부수성이 존재할 수 있다고 주장했습니다 --- 아리스토텔레스 자신의 정의(\citework{Cat}{1a24--25})와 정면 충돌합니다.

1277년 파리 금지령에서 제140조(부수성의 주어 없는 존재가 불가능하다는 명제)를 이단으로 규정함으로써, 아리스토텔레스 형이상학의 자율성이 제한되었습니다.\footnote{Edward Grant, ``The Effect of the Condemnation of 1277,'' in \gri{The Cambridge History of Later Medieval Philosophy}, eds.\ N.\ Kretzmann, A.\ Kenny, J.\ Pinborg, ch.\ 26 (Cambridge: Cambridge University Press, 1982).}

\newpage
\subsection{돌아보기와 앞으로의 길}

\subsubsection{세 문명권의 비교 --- 같은 텍스트, 다른 물음}

같은 두 텍스트가 서로 다른 핵심 물음을 촉발했습니다.

{\footnotesize
\begin{concepttable}{|l|l|l|}
\hline
\rowcolor{tableheadblue}
{\color{h1navy}\bfseries 문명권} &
{\color{h1navy}\bfseries 핵심 물음} &
{\color{h1navy}\bfseries 대표적 논쟁} \\ \hline
그리스 & 범주론은 무엇에 관한 책인가? & 말/개념/사물 \\ \hline
이슬람 & 범주론은 논리학에 속하는가? & 1차/2차 지향 \\ \hline
비잔틴 & 유산을 어떻게 보존하고 신학과 조화시키는가? & 이탈로스 정죄 (1082) \\ \hline
라틴 서방 & 보편자는 존재하는가? & 실재론/유명론 \\ \hline
(신학) & 실체 없는 부수성은 가능한가? & 화체론/1277년 금지령 \\ \hline
\end{concepttable}
}

\subsubsection{범주론에서 명제론으로}

범주론은 비결합 표현들을 분석했습니다. 이사고게는 술어화 방식을 분석했습니다. 다음 주(9주차)부터 시작하는 명제론은 이 표현들이 결합되어 명제를 이루었을 때 참이거나 거짓이 되는 구조를 탐구합니다.

명제론의 첫 줄(16a1)은 의미의 삼각형으로 시작합니다 --- 발화된 말은 영혼의 경험들의 상징(\gr{σύμβολα})이고, 영혼의 경험들은 사물들의 모상(\gr{ὁμοιώματα})이다. 이것은 특별 세션 α에서 다룬 심플리키오스의 `원래 하나였던 세 가지'(말, 개념, 사물)를 아리스토텔레스 자신이 서술하는 대목이며, 범주론과 이사고게에서 우리가 분석한 모든 것의 출발점이 됩니다.

\subsubsection{우리가 배운 읽기의 기술}

7주간의 강독과 2주간의 조망 세션을 통해 우리는 내용뿐 아니라 읽기의 기술도 배웠습니다. 그리스어 원문과 한국어 번역을 병행하며 읽는 것, 2차 문헌과 대화하며 읽는 것, 텍스트 내부의 긴장을 발견하고 열린 문제로 유지하는 것, 수용사를 통해 한 텍스트가 여러 문명에서 어떻게 다르게 읽혔는지를 추적하는 것 --- 이것은 남은 34주 동안 분석론, 토피카, 소피스트적 논박을 읽는 데 계속 사용될 기술입니다.

\subsection*{참고문헌}
\begin{references}

\refitem Benakis, K. ``Michael Psellos.'' In \gri{Philosophie: Geschichtliche Einführungen}, vol.\ 6. Basel, 1982.

\refitem Cerami, Cristina. ``Avicenna against Porphyry's Definition of Differentia Specifica.'' \gri{Documenti e studi sulla tradizione filosofica medievale} 26 (2015): 129--183.

\refitem De Rijk, L.\ M. \gri{Logica Modernorum}. Assen: Van Gorcum, 1962.

\refitem Grant, Edward. ``The Effect of the Condemnation of 1277.'' In \gri{The Cambridge History of Later Medieval Philosophy}, edited by N.\ Kretzmann, A.\ Kenny, J.\ Pinborg, ch.\ 26. Cambridge: Cambridge University Press, 1982.

\refitem Hodges, Wilfrid. ``al-Farabi's Philosophy of Logic and Language.'' In \gri{Stanford Encyclopedia of Philosophy}, 2023.

\refitem Ierodiakonou, Katerina. ``Logic in Byzantium.'' In \gri{The Cambridge Intellectual History of Byzantium}, edited by A.\ Kaldellis and N.\ Siniossoglou. Cambridge: Cambridge University Press, 2017.

\refitem Ighbariah, Ahmad. ``Al-Kind\=\i{} and the Reception of Aristotle's Categories.'' \gri{Arabic Sciences and Philosophy} 22, no.\ 1 (2012).

\refitem Klima, Gyula. ``The Medieval Problem of Universals.'' In \gri{Stanford Encyclopedia of Philosophy}, rev.\ 2022.

\refitem Spade, Paul Vincent. \gri{Five Texts on the Mediaeval Problem of Universals}. Indianapolis: Hackett, 1994.

\refitem Strobino, Riccardo. ``Ibn Sina's Logic.'' In \gri{Stanford Encyclopedia of Philosophy}.

\end{references}

\end{document}
